% Part: turing-machines 
% Chapter: undecidability
% Section: introduction

\documentclass[../../../include/open-logic-section]{subfiles}

\begin{document}

\olfileid{tur}{und}{int} 
\olsection{प्रस्तावना}

प्रत्येक फलन, अगदी प्रत्येक अंकगणितीय फलनही, संगणनक्षम
असेलच असे नाही, हे उघड वाटू शकते. अशी फलने खूप आहेत
आणि त्यांचे वर्तन फार गुंतागुंतीचे असू शकते. उदाहरणार्थ,
किरणोत्सर्गी कणांच्या क्षयावरून किंवा इतर अव्यवस्थित अथवा
यादृच्छिक वर्तनावरून परिभाषित केलेली फलने. समजा, एखाद्या
दिलेल्या क्षणापासून आपण एका सेकंदाच्या अंतरांची मोजणी सुरू
करतो आणि त्या आरंभीच्या क्षणानंतरच्या $n$-व्या एक-सेकंदाच्या
अंतरात विश्वातील जितक्या कणांचा क्षय होतो ती संख्या $f(n)$
या फलनाची किंमत म्हणून ठरवतो. हे फलन कधीही संगणित
करता येणार नाही असे वाटण्यासारखे आहे.

मात्र अशी फलने कशी संगणित करावीत याची कल्पना न सुचणे
ही एक गोष्ट आहे; ती संगणनक्षम नाहीत हे प्रत्यक्ष सिद्ध
करणे ही दुसरी गोष्ट आहे. खरे तर, अत्यंत गुंतागुंतीची
वाटणारी फलनेही अमूर्त अर्थाने संगणनक्षम असू शकतात.
उदाहरणार्थ, विश्वाचा कालावधी मर्यादित आहे असे समजा---काही
विश्वशास्त्रीय सिद्धांतांच्या भाकिताप्रमाणे फार दूरच्या भविष्यात
विश्व आकुंचन पावून एका बिंदूत सामावेल. मग त्या आरंभीच्या
क्षणापासून $f(n)$ परिभाषित असेल असे फक्त मर्यादित
(पण अविश्वसनीयरीत्या मोठ्या) संख्येने सेकंद असतील.
आणि केवळ मर्यादित संख्येच्या आदानांवर परिभाषित असलेले
कोणतेही फलन संगणनक्षम असते: त्याची फलिते एका मोठ्या
सारणीमध्ये नोंदवता येतील किंवा ट्यूरिंग यंत्राच्या एका
अतिशय मोठ्या अवस्था-संक्रमण आलेखात सांकेतिक करता येतील.

ज्या फलनांच्या किंमती होय/नाही या प्रश्नांची उत्तरे देतात,
अशा फलनांच्या विशेष प्रकारांत आपल्याला अनेकदा रस असतो.
उदाहरणार्थ, “$n$ ही अविभाज्य संख्या आहे का?” हा प्रश्न
पुढील फलनाशी संबंधित आहे:
\[
\fn{isprime}(n) = \begin{cases}
  1 & \text{जर\/ $n$ अविभाज्य असेल तर}\\
  0 & \text{अन्यथा.}
  \end{cases}
\]
एखाद्या होय/नाही प्रश्नाशी संबंधित $1/0$-मूल्यी फलन
प्रभावी रीतीने संगणनक्षम असेल, तर त्या प्रश्नाचा
\emph{प्रभावी रीतीने निर्णय करता येतो} असे आपण म्हणतो.

प्रभावी रीतीने संगणित करता न येणारी फलने किंवा प्रभावी
रीतीने निर्णय करता न येणाऱ्या समस्या अस्तित्वात आहेत,
हे गणिती पद्धतीने सिद्ध करण्यासाठी संगणनाचे विशिष्ट
प्रतिमान निश्चित करणे आवश्यक आहे. त्यानंतर त्या
प्रतिमानाला काही फलने संगणित करता येत नाहीत किंवा
काही समस्यांचा निर्णय करता येत नाही, हे दाखवावे लागते.
उदाहरणार्थ, सर्वच फलने ट्यूरिंग यंत्रांनी संगणित करता
येत नाहीत आणि सर्वच समस्यांचा निर्णय ट्यूरिंग यंत्रांनी
करता येत नाही, हे आपण दाखवू शकतो. मग चर्च--ट्यूरिंग
प्रबंधाच्या आधारे केवळ ट्यूरिंग यंत्रे सर्व फलने संगणित
करण्याइतकी सामर्थ्यवान नाहीत असेच नव्हे, तर कोणतीही
प्रभावी प्रक्रियाही तसे करू शकत नाही, असा निष्कर्ष
काढता येतो.

असे निषेधात्मक निष्कर्ष सिद्ध करण्याची गुरुकिल्ली म्हणजे
ट्यूरिंग यंत्रांनाही संख्या देता येतात ही गोष्ट. याचा
सर्वांत सोपा मार्ग म्हणजे त्यांचे प्रगणन करणे: ट्यूरिंग
यंत्रे आणि त्यांचे कार्यक्रम लिहिण्याची एखादी निश्चित
पद्धत ठरवून त्यांची पद्धतशीर यादी करता येते. असे
करता येते हे लक्षात आल्यावर, केवळ संचसंख्येच्या
विचारांवरून ट्यूरिंग-असंगणनक्षम फलने अस्तित्वात आहेत
हे कळते: $\Nat$ पासून~$\Nat$ मध्ये जाणाऱ्या फलनांचा
संच (खरे तर, $\Nat$ पासून $\{0, 1\}$ मध्ये जाणाऱ्या
फलनांचा संचसुद्धा) !!{nonenumerable} म्हणजे अगणनीय
आहे; परंतु सर्व ट्यूरिंग यंत्रांचे प्रगणन करता येत
असल्यामुळे ट्यूरिंग-संगणनक्षम फलनांचा संच केवळ
!!{denumerable} म्हणजे गणनीय अनंत आहे.

याशिवाय, संगणित करता येत नाहीत किंवा ज्यांचा निर्णय
करता येत नाही अशा अनुक्रमे \emph{विशिष्ट} फलनांची
आणि समस्यांचीही व्याख्या करता येते. त्यांपैकी
एक म्हणजे तथाकथित \emph{थांबण्याची समस्या.}
ट्यूरिंग यंत्रांच्या सूचना यादीत दिल्यास त्यांचे
मर्यादित वर्णन करता येते. ट्यूरिंग यंत्राचे असे
वर्णन, म्हणजेच ट्यूरिंग यंत्राचा कार्यक्रम, अर्थातच
दुसऱ्या ट्यूरिंग यंत्राला आदान म्हणून देता येते.
म्हणून इतर ट्यूरिंग यंत्रांबद्दलच्या प्रश्नांचा
निर्णय करणाऱ्या ट्यूरिंग यंत्रांचा विचार करता येतो.
विशेष रोचक प्रश्न असा आहे: “दिलेल्या ट्यूरिंग यंत्राला
आदान~$n$ दिल्यावर ते अखेरीस थांबते का?”
या प्रश्नाचा निर्णय करू शकणारे ट्यूरिंग यंत्र असते
तर बरे झाले असते: ट्यूरिंग यंत्रे अनंत चक्रात अडकत
नाहीत वगैरे तपासणारे गुणवत्ता-नियंत्रण यंत्र म्हणून
त्याचा विचार करा. ट्यूरिंग यांनी सिद्ध केलेली
रोचक गोष्ट म्हणजे असे यंत्र असूच शकत नाही.
ट्यूरिंग यंत्र~$M$ चे वर्णन आणि एखादी
संख्या~$n$ यांचे आदान दिल्यावर $M$ हे यंत्र
आदान $n$ वर थांबले असते की नाही यानुसार नेहमी
थांबून अनुक्रमे $1$ किंवा $0$ फलित देणारे एकही
ट्यूरिंग यंत्र असू शकत नाही.

अनिर्णेय समस्यांची विशिष्ट उदाहरणे मिळाल्यावर इतर
समस्याही अनिर्णेय आहेत हे दाखवण्यासाठी त्यांचा
उपयोग करता येतो. उदाहरणार्थ, “प्रथम-क्रम
!!{formula}~$!A$ वैध आहे का?” हा एक प्रसिद्ध
अनिर्णेय प्रश्न आहे. आदान म्हणून प्रथम-क्रम
!!{formula}~$!A$ दिल्यावर $!A$ वैध आहे की
नाही यानुसार $1$ किंवा $0$ फलित देत नेहमी
थांबेल असे कोणतेही ट्यूरिंग यंत्र नाही.
ऐतिहासिकदृष्ट्या, ही समस्या प्रभावी रीतीने
सोडवण्याची प्रक्रिया शोधण्याच्या प्रश्नाला
सरळ “निर्णय समस्या” म्हटले जात असे; म्हणून
निर्णय समस्या सोडवता येत नाही असे आपण म्हणतो.
ट्यूरिंग आणि चर्च यांनी हे फलित साधारणतः
एकाच काळात स्वतंत्रपणे सिद्ध केले; त्यामुळे
त्याला चर्च--ट्यूरिंग प्रमेय असेही म्हणतात.

\end{document}
